Artin's braid group $B_n$ appears throughout mathematics in many guises, in large part due to its fundamental connection to the motion of sets of particles in the plane. The group $B_n$ enjoys many desirable properties, such as finite-presentability and linearity. The latter was established via the faithful \emph{Lawrence--Krammer representation} of $B_n$; however, a long-standing open problem is for which values of $n$ the related \emph{Burau representation} $\Psi_n$ of the braid group is faithful. It is known to be faithful for $n \leq 3$ and not faithful for $n \geq 5$. In this paper, we are concerned with the question of faithfulness when $n=4$, the remaining open case. A non-trivial element in the kernel of $\Psi_4$ would give a likely candidate for showing that the Jones polynomial cannot detect the unknot, as noted by Bigelow \cite{Big02}.

The Burau representation's lack of faithfulness for large $n$ was shown by Moody \cite{Moo91}, Long--Paton \cite{LP93}, and Bigelow \cite{Big99}, using related criteria involving combinatorics of pairs of arcs in the $n$-punctured disk $D_n$. To each pair of arcs in $D_n$ we associate a polynomial called its \emph{Burau polynomial}. One criterion states that $\Psi_n$ is not faithful if and only if there exists a pair of intersecting arcs in an explicit set $\mathcal{A}$ whose Burau polynomial is zero. Ordering arc-pairs by their geometric intersection number, the first main theorem of this paper certifies the faithfulness of the Burau representation $\Psi_4$ up to arc-pairs with intersection number 2000.

\begin{maintheorem}\label{the:certify}Any pair of arcs $(\alpha, \beta) \in \mathcal{A}$ with geometric intersection number at most 2000 has non-zero Burau polynomial.
\end{maintheorem} 

We prove Theorem~\ref{the:certify} by computer search, identifying then testing a finite yet sufficient collection of arc-pair representatives for each intersection number. We describe the arcs using weighted train tracks, with linear inequalities between the weights being used to `steer' the arcs around the disk $D_n$.

Computation of Burau polynomials has been carried out previously. Working with a different collection of arcs in $D_n$ referred to as `noodles' and `forks', Bigelow \cite{Big02} certified the faithfulness of $\Psi_4$ up to noodles and forks intersecting 2000 times. Translating this into the language used here, this corresponds to a set of arc-pairs intersecting at most 1000 times. The bound obtained in Theorem~\ref{the:certify} is thus of notable increase. Also, the collection of arcs in \cite{Big02} and those in Theorem~\ref{the:certify} are of entirely different types. While either collection may be used to certify faithfulness, we are unaware of any way to compare the two.  Our approach benefited from decreased computation time, due to several short-cuts we were able to make. These are discussed in detail in Section 4.

One by-product of the proof of Theorem~\ref{the:certify} is the means to list a slew of Burau polynomials. In doing so, a striking periodicity is observed, as we now describe. We define the \emph{norm} of a Burau polynomial to be the sum of the absolute values of its coefficients. For reasons apparent in the calculation of these polynomials, if seeking the zero polynomial, we need only check arc-pairs with even intersection number. The following theorem records a repetitive structure observed in the minimum norm for such intersection numbers.

\begin{maintheorem}\label{the:period}Let $2k \in [44, 500]$. If $2k$ is divisible by $6$, then the minimum norm over all Burau polynomials arising from an arc-pair with $2k$ geometric intersections is $10$. Otherwise, the minimum norm is $8$.\end{maintheorem}

The upper bound of 500 was imposed by limits on computation time; this periodicity almost certainly persists longer. Immediately, we observe that if $\Psi_4$ is not faithful, there must be some large intersection number at which the periodicity breaks.

We remark that the largest intersection numbers checked in Theorems \ref{the:certify} and \ref{the:period} are of different magnitudes. Both were pushed as high as possible, using hundreds of hours of computation time. However, the former only required testing whether the given polynomials equaled zero, whereas the latter demanded that we form a complete list of Burau polynomials to examine. For a given geometric intersection number, the latter case is more computationally laborious.

A strong relationship is observed between the arcs of minimum norm that appear in Theorem~\ref{the:period}. Every such arc falls into a family, where within a single family, the train track weights lie in an arithmetic progression and the ordered sets of coefficients of their Burau polynomials are all the same. Attempting to see the behavior of Burau polynomials across all arc-pairs, these `minimum norm families' obscure our view. By grouping all arcs into similar families, we are able to see through the minimum norms obtained for each intersection number and give numerical evidence that suggests $\Psi_4$ is faithful. This is discussed further in Section 5.

\textbf{Outline of paper.} In Section 2, we recall the definition of the Burau representation, and state criteria for its faithfulness. Section 3 explains how we use weighted train tracks to describe arcs in the disk $D_n$, and how we are able to reduce to only checking arc-pairs of a certain form. Section 4 details the algorithm we constructed, and in Section 5 we report on our observations and numerical evidence.

\textbf{Acknowledgements.} The authors would like to thank Stephen Bigelow, Dan Margalit, and Bal\'azs Strenner for helpful conversations. Part of this work was completed while the first author was in residence at the Mathematical Sciences Research Institute in Berkeley, California, during the Fall 2016 semester.